\begin{abstract}
Deep learning models for visual recognition often exhibit systematic errors due to underrepresented semantic subpopulations. Although existing debugging frameworks can pinpoint these failures by identifying key failure attributes, repairing the model effectively remains difficult. Current solutions often rely on manually designed prompts to generate synthetic training images—an approach prone to distribution shift and semantic errors. To overcome these challenges, we introduce a model repair module that builds on an interpretable failure attribution pipeline. Our approach uses a conditional text-to-image model to generate semantically faithful and targeted images for failure cases. To preserve the quality and relevance of the generated samples, we further employ a large vision--language model (LVLM) to filter the outputs, enforcing alignment with the original data distribution and maintaining semantic consistency. By retraining vision models with this rare-case-augmented synthetic dataset, we significantly reduce errors associated with rare cases. Our experiments demonstrate that this targeted repair strategy improves model robustness without introducing new bugs.
\end{abstract}

\section{Introduction}
Despite extensive training on annotated images and evaluation on static test sets, computer vision models can still fail in unexpected ways~\cite{barbu2019objectnet,gao2023adaptive,leclerc20223db}. One common cause is dataset bias: certain subpopulations are underrepresented in the training data. For instance, facial recognition systems have shown significantly higher error rates for darker-skinned women compared to lighter-skinned men, due to imbalanced representation in training datasets~\cite{buolamwini2018gender}. Therefore, identifying and repairing potential model failures caused by underrepresented semantic subpopulations before deployment is essential to improve the reliability of computer vision systems in real-world applications.

Several recent studies have addressed the identification and repair of model failures caused by underrepresented semantic subpopulations. For example, given a natural language description of a potentially underrepresented semantic subpopulation, AdaVision~\cite{gao2023adaptive} introduces a model debugging pipeline that leverages CLIP~\cite{radford2021learning} to iteratively retrieve relevant images from LAION-5B~\cite{schuhmann2022laion} and evaluate model predictions. \safefix~\cite{chen2023safefix} advances interpretable model debugging by automatically discovering failure slices based on shared visual attributes using a pre-trained vision--language model. DCD~\cite{singla2024data} identifies failure modes by collecting misclassified samples, computing a failure direction for distinguishing attributes, and retrieving similar unlabeled examples to expose systematic weaknesses and underrepresented subpopulations.
The latent-space method~\cite{jain2023latentfailures} discovers failure patterns by training a linear SVM in a  latent space, separating correct from incorrect predictions to yield subpopulation errors without labels.
D3M~\cite{jain2023d3m} enhances overall model performance by removing harmful training samples that disproportionately affect underrepresented groups.

While existing methods can effectively identify model errors caused by underrepresented semantic subpopulations, the subsequent and critical step—model repair—still presents challenges. For instance, AdaVision~\cite{gao2023adaptive} employs a simple fine-tuning strategy to retrain the model on retrieved error-causing images.
However, this approach can degrade performance on images where the model originally performed well. \safefix~\cite{chen2023safefix} addresses model repair by turning attributes describing underrepresented subpopulations into language prompts to synthesize new training images via generative models. However, the synthesized images may suffer from \emph{distribution mismatch} compared to the original training data. Moreover, the repair process can introduce new bugs, as the generated images may fail to preserve attributes unrelated to original model errors.


Building on these observations, we propose \model, a targeted model repair method designed to correct failures caused by underrepresented semantic subpopulations through controlled image generation—without degrading overall performance or introducing new errors. Compared to AdaVision~\cite{gao2023adaptive} and DCD~\cite{singla2024data}, which retrieve samples from another dataset, our method generates new examples grounded in the training set to ensure distributional consistency. Unlike \safefix~\cite{chen2023safefix}, which simply composes failure attributes into language prompts, \model anchors generation on real instances in the training set to preserve unrelated attributes and avoid unintended changes. While the latent-space-based filtering method~\cite{jain2023latentfailures} explores repair via diffusion models, their approach relies solely on caption-based prompting and does not guarantee the semantic accuracy of critical attributes. In contrast, \model mitigates the unreliability of generative models in rendering sensitive attributes—such as skin tone or hair color—by further filtering out semantically inconsistent images using a large vision--language model (LVLM)~\cite{bai2025qwen25vl, liu2023improvedllava}, and the LVLM is verified by human audit. This ensures that the final outputs faithfully reflect the intended attribute change (e.g., generating a sad expression on a darker-skinned woman with red hair, whereas standard diffusion models often mistakenly produce a darker-skinned man with black hair instead), enabling precise and robust repair of underrepresented semantic subpopulation failures. 


\model directly addresses two major challenges in existing model repair strategies: {\bf 1)} ensuring the generated images reflect the true failure attributes, and {\bf 2)} aligning them with the original data distribution. By retraining models on this augmentation set, we reduce errors associated with underrepresented semantic subpopulations, where model failures stem from insufficient training coverage of specific attributes. Our contributions are as follows:
\begin{itemize}
\item We propose a targeted model repair pipeline, \model, which leverages conditional text-to-image generation and LVLM-based filtering to synthesize high-quality, attribute-faithful data for correcting model failures arising from underrepresented subpopulations.

\item We show that our approach mitigates distribution mismatch in generated samples, ensuring that the augmented data remains aligned with the original dataset distribution. This alignment helps correct model failures while maintaining overall performance.
\item We demonstrate that retraining with curated, distribution-consistent synthetic data significantly enhances model robustness in underrepresented scenarios—without introducing new errors.

\end{itemize}